\documentclass[11pt,a4paper]{article}
\input{preamble}
\title{Principles of Quantum Mechanics}
\begin{document}
\maketitle
\thispagestyle{empty}
\section*{0. Introduction}

\noindent
$\ast$ Recall features from IB QM (Supplementary Notes 1)

\begin{itemize}
    \item Wave-particle duality:
    
    Waves behaving like particles, e.g.\ light quanta,
    \& vice versa (particles behaving like waves).
    
    \item Wavefunction $\psi(x)$ for a particle --- prob.\ density
    $|\psi(x)|^2$.
    
    (Prob.\ intrinsic to theory.)
    
    \item Observables became operators on wavefns:
    \[
        \hat{x} \to x,
        \qquad
        \hat{p} \to -i\hbar\frac{\partial}{\partial x}.
    \]
    \& lack of commutativity, e.g.
    \[
        [\hat{x},\hat{p}] = i\hbar
    \]
    leads to Uncertainty Relation.
    
    \item Schr\"odinger eqn with Hamiltonian $H$ specifies dynamics
    \& eigenvalues of $H$ are allowed energy levels of system.
    
    \item All of this explains H atom structure \& resulting line spectra.
\end{itemize}

\noindent
$\ast$ Aims of this course:

\begin{itemize}
    \item Reformulate QM in a more powerful, flexible, \& useful form
    --- the Dirac formalism (go beyond wavefns!).
\end{itemize}

Previous description appears as particular realisation
\& also leads to alternatives, e.g.
\[
    \tilde{\psi}(p)
    \qquad \text{with} \qquad
    \hat{x} \to +i\hbar\frac{\partial}{\partial p}
    \quad \& \quad
    \hat{p} \to p.
\]

This more general approach allows
\begin{itemize}
    \item Simpler analysis of known problems, e.g.\ oscillator
\end{itemize}

but also lots of new things
\begin{itemize}
    \item \underline{Spin} of particles
    
    \item \underline{Symmetries} (translations, rotations, for example)
    \& conservation laws.
    
    \item \underline{Identical particles} (bosons \& fermions)
    
    \item Framework for quantising more general systems,
    e.g.\ EM field
    
    \item \underline{Perturbation theory}: practical tool for computations ---
    we will not look at too many detailed applications, but we
    will try to keep track of what all this is for --- physical predictions
\end{itemize}

\section*{1. Dirac Formalism}

\subsection*{1.1 Concepts \& Definitions}

A quantum system is described at each instant by a \underline{state}
$|\psi\rangle$ which belongs to a complex vector space $V$;
any non-zero element of this space $V$ is a possible state.

Thus
\[
    |\psi\rangle \ \&\ |\phi\rangle \in V
    \quad\Rightarrow\quad
    \alpha|\psi\rangle+\beta|\phi\rangle\in V
    \qquad
    \text{for any }\alpha,\beta\in\mathbb{C}.
\]

Physically this is a superposition principle leading to wave-like behaviour,
but these states are \underline{not} \underline{wavefns}.
There are also \underline{dual} or conjugate states
$\langle\phi|$ which belong to the dual space $V^\dagger$.

By definition, states \& duals can be combined or paired to give a
complex number
\[
    \underset{\text{``bra''}}{\langle\phi|},
    \quad
    \underset{\text{``ket''}}{|\psi\rangle}
    \ \longmapsto\
    \underset{\text{``bra(c)ket''}}{\langle\phi|\psi\rangle},
    \qquad
    V^\dagger\times V\to\mathbb{C},
\]
with
\begin{align*}
    \langle\phi|
    \bigl(\alpha_1|\psi_1\rangle+\alpha_2|\psi_2\rangle\bigr)
    &=
    \alpha_1\langle\phi|\psi_1\rangle
    +\alpha_2\langle\phi|\psi_2\rangle,\\
    \bigl(\beta_1\langle\phi_1|+\beta_2\langle\phi_2|\bigr)
    |\psi\rangle
    &=
    \beta_1\langle\phi_1|\psi\rangle
    +\beta_2\langle\phi_2|\psi\rangle.
\end{align*}

The spaces $V$ \& $V^\dagger$ come with an
\underline{inner product}, which can be described as a one-to-one
correspondence between the spaces
\[
    V\longleftrightarrow V^\dagger,
\]
with
\[
    |\psi\rangle
    \longleftrightarrow
    \langle\psi|
    =
    \bigl(|\psi\rangle\bigr)^\dagger
    \qquad
    \text{(use same symbol for corresponding states).}
\]

The inner product is
\[
    |\phi\rangle,\ |\psi\rangle
    \longmapsto
    \langle\phi|\psi\rangle
    =
    \bigl(|\phi\rangle\bigr)^\dagger|\psi\rangle,
    \qquad
    V\times V\to\mathbb{C},
\]
\& it is assumed to obey
\[
    \langle\phi|\psi\rangle
    =
    \langle\psi|\phi\rangle^*
    \qquad
    \underline{\text{hermitian}},
\]
\[
    \bigl\|\,|\psi\rangle\,\bigr\|^2
    =
    \langle\psi|\psi\rangle
    \qquad
    \text{real},
\]
\[
    \left.
    \begin{aligned}
        \langle\psi|\psi\rangle &\geq 0,\\
        \langle\psi|\psi\rangle &=0
        \quad\text{iff}\quad
        |\psi\rangle=0
    \end{aligned}
    \right\}
    \qquad
    \underline{\text{positive definite}},
\]
where now $\bigl\|\,|\psi\rangle\,\bigr\|$ is called the norm.

\medskip

\noindent
Note: Knowing $\langle\phi|\psi\rangle$ for all
$\langle\phi|$ determines $|\psi\rangle$
(\& vice versa).

\medskip

\noindent
[Check:
\[
    \langle\phi|\psi_1\rangle
    =
    \langle\phi|\psi_2\rangle
    \quad\Rightarrow\quad
    \langle\phi|\chi\rangle=0,
    \qquad
    \text{where }
    |\chi\rangle=|\psi_1\rangle-|\psi_2\rangle.
\]
If this holds $\forall\,\langle\phi|$, set
$\langle\phi|=\langle\chi|$ \& get
\[
    \langle\chi|\chi\rangle
    =
    \bigl\|\,|\chi\rangle\,\bigr\|^2
    =
    0
    \quad\Rightarrow\quad
    |\chi\rangle=0,
    \qquad
    \text{as claimed.}
\]
]

\medskip

The physical content of any state is unchanged by
\[
    |\psi\rangle\to\alpha|\psi\rangle
    \qquad
    \text{for any }\alpha\in\mathbb{C},\quad\alpha\neq0.
\]

We will usually normalise to make
$\bigl\|\,|\psi\rangle\,\bigr\|=1$, but can still change
\[
    |\psi\rangle\to e^{i\theta}|\psi\rangle.
\]

Relative phases can be important in expressions like
$\alpha|\psi\rangle+\beta|\phi\rangle$.

The space $V$ is \underline{complete}: we assume that appropriate
sequences \& series \underline{converge}.
A complete inner product space is called a
\underline{Hilbert space}.

\medskip

An \underline{operator} $Q$ is a linear map on states
\[
    |\psi\rangle\longmapsto Q|\psi\rangle,
    \qquad
    V\to V,
\]
with
\[
    Q\bigl(\alpha|\psi\rangle+\beta|\phi\rangle\bigr)
    =
    \alpha Q|\psi\rangle+\beta Q|\phi\rangle.
\]

The same operator can be thought of as acting to the left on dual states,
\[
    \langle\phi|\longmapsto\langle\phi|Q,
    \qquad
    V^\dagger\to V^\dagger,
\]
where the right-hand side is defined by
\[
    \bigl(\langle\phi|Q\bigr)|\psi\rangle
    =
    \langle\phi|\bigl(Q|\psi\rangle\bigr)
    =
    \langle\phi|Q|\psi\rangle.
\]

\begin{center}
    \textbf{End of lecture 1}
\end{center}
Last time: abstract notation for states
\[
    |\psi\rangle\in V,
    \qquad
    \langle\phi|\in V^\dagger,
    \qquad
    \text{\& map }V\longleftrightarrow V^\dagger,
\]
\& operators acting, e.g.\ $|\psi\rangle\longmapsto Q|\psi\rangle$
(\S\,1.1).

\medskip

\noindent
Note: States \& operators $|\psi\rangle$, $\langle\phi|$, $Q$
can be dealt with by rules that mirror those for complex vectors
$\underline{u}$ (column), $\underline{v}^{\dagger}$ (row),
\& a matrix $A$; e.g.\ compare
\[
    \langle\phi|Q|\psi\rangle
    \qquad\text{with}\qquad
    \underline{v}^{\dagger}A\underline{u}.
\]

For any operator $Q$, say $|\psi\rangle\neq0$ is an
\underline{eigenstate} with \underline{eigenvalue} $\lambda$ if
\[
    Q|\psi\rangle=\lambda|\psi\rangle;
\]
equivalently,
\[
    \langle\psi|Q^\dagger=\lambda^*\langle\psi|.
\]
For general $Q$, $\lambda$ may be complex.

For any $Q$, the \underline{hermitian conjugate} or adjoint
$Q^\dagger$ is defined by
\[
    \langle\phi|Q^\dagger
    =\bigl(Q|\phi\rangle\bigr)^\dagger
    \qquad\text{for any }|\phi\rangle.
\]
Equivalently,
\begin{align*}
    \langle\phi|Q^\dagger|\psi\rangle
    &=\bigl(Q|\phi\rangle\bigr)^\dagger|\psi\rangle\\
    &=\langle\psi|Q|\phi\rangle^*
    \qquad\text{for all }|\phi\rangle,\ |\psi\rangle.
\end{align*}
Simple consequences:
\begin{align*}
    (\alpha A+\beta B)^\dagger
    &=\alpha^*A^\dagger+\beta^*B^\dagger,\\
    (AB)^\dagger&=B^\dagger A^\dagger,
\end{align*}
for any operators $A,B$ \& $\alpha,\beta\in\mathbb{C}$, \&
\[
    \bigl(Q^\dagger\bigr)^\dagger=Q.
\]

\medskip

\noindent
\underline{Observables}

An operator $Q$ is \underline{hermitian} or
\underline{self-adjoint} if
\[
    Q^\dagger=Q;
\]
such operators are called observables because they correspond to
physical quantities that can be measured, e.g.\ position, momentum, energy.

To define a quantum system, we must specify the \underline{Hamiltonian}
$H$ whose eigenvalues are the allowed energy levels of the system:
\[
    H|\chi\rangle=E|\chi\rangle
    \qquad\text{energy eigenvalue problem.}
\]

Key results for any hermitian $Q$:
\begin{enumerate}
    \item[(i)] All eigenvalues are real.
    \item[(ii)] Eigenstates with distinct eigenvalues are orthogonal.
    \item[(iii)] Any state can be expanded in terms of
    (written as a linear combination of) eigenstates.
\end{enumerate}
We prove (i) \& (ii) \& assume (iii).

\medskip

\noindent
(i) Suppose
\[
    Q|\psi\rangle=\lambda|\psi\rangle.
\]
Then
\[
    \langle\psi|Q^\dagger
    =\lambda^*\langle\psi|
    =\langle\psi|Q
    \qquad(Q^\dagger=Q).
\]
Hence
\[
    \langle\psi|Q|\psi\rangle
    =\lambda\langle\psi|\psi\rangle
    =\lambda^*\langle\psi|\psi\rangle.
\]
But
\[
    \bigl\|\,|\psi\rangle\,\bigr\|^2
    =\langle\psi|\psi\rangle\neq0
    \qquad(|\psi\rangle\neq0),
\]
so $\underline{\lambda^*=\lambda}$.

\medskip

\noindent
(ii) Let $|n\rangle$ be eigenstates of $Q$ with eigenvalues
$\lambda=q_n$ (real), where $n$ is some discrete label,
possibly ranging over infinitely many values.
\[
    Q|n\rangle=q_n|n\rangle,
    \qquad
    Q|m\rangle=q_m|m\rangle,
\]
or
\[
    \langle m|Q=q_m\langle m|
    \qquad(q_m^*=q_m).
\]
Hence
\[
    \langle m|Q|n\rangle
    =q_n\langle m|n\rangle
    =q_m\langle m|n\rangle.
\]
But $q_m\neq q_n$, so $\langle m|n\rangle=0$.

\subsection*{1.2 The Harmonic Oscillator}

A one-dimensional harmonic oscillator of mass $m$ \& frequency
$\omega$ is a system with position $\hat{x}$, momentum $\hat{p}$
\& Hamiltonian
\[
    H=\frac{1}{2m}\hat{p}^{\,2}
    +\frac{1}{2}m\omega^2\hat{x}^{\,2}.
\]
Seek eigenstates \& eigenvalues of $H$.

\medskip

\noindent
(a) \underline{Annihilation, creation \& number operators}

Define
\[
    \left\{
    \begin{aligned}
        a&=\left(\frac{m\omega}{2\hbar}\right)^{1/2}
        \left(\hat{x}+\frac{i\hat{p}}{m\omega}\right),\\
        a^\dagger&=\left(\frac{m\omega}{2\hbar}\right)^{1/2}
        \left(\hat{x}-\frac{i\hat{p}}{m\omega}\right).
    \end{aligned}
    \right.
\]
These are dimensionless hermitian conjugates, since
$\hat{x}^\dagger=\hat{x}$, $\hat{p}^\dagger=\hat{p}$.
Equivalently,
\[
    \left\{
    \begin{aligned}
        \hat{x}&=\left(\frac{\hbar}{2m\omega}\right)^{1/2}
        (a+a^\dagger),\\
        \hat{p}&=\left(\frac{\hbar m\omega}{2}\right)^{1/2}
        i(a^\dagger-a).
    \end{aligned}
    \right.
\]
Easy to check:
\[
    [\hat{x},\hat{p}]=i\hbar
    \quad\Longleftrightarrow\quad
    [a,a^\dagger]=1.
\]
In fact,
\begin{align*}
    aa^\dagger
    &=\frac{m\omega}{2\hbar}\hat{x}^{\,2}
    +\frac{1}{2m\omega\hbar}\hat{p}^{\,2}
    -\frac{i}{2\hbar}(\hat{x}\hat{p}-\hat{p}\hat{x})\\
    &=\frac{1}{\hbar\omega}H+\frac{1}{2},
\end{align*}
but
\[
    a^\dagger a=\frac{1}{\hbar\omega}H-\frac{1}{2}.
\]
This confirms the commutation relation, but we also find
\[
    H=\hbar\omega\left(N+\frac{1}{2}\right),
\]
with $N=a^\dagger a$ the \underline{number operator}.

$N^\dagger=N$, \& finding eigenvalues \& eigenstates of $N$
is equivalent to finding those for $H$.
Let $|\lambda\rangle$ be any eigenstate of $N$ with eigenvalue $\lambda$:
\[
    N|\lambda\rangle=\lambda|\lambda\rangle
    \quad\Longleftrightarrow\quad
    H|\lambda\rangle=E|\lambda\rangle,
    \qquad
    E=\hbar\omega\left(\lambda+\frac{1}{2}\right).
\]
Then, assuming $\bigl\|\,|\lambda\rangle\,\bigr\|^2=1$
(normalised),
\begin{align*}
    \lambda
    &=\langle\lambda|N|\lambda\rangle
    =\langle\lambda|a^\dagger a|\lambda\rangle
    =\bigl(a|\lambda\rangle\bigr)^\dagger
    \bigl(a|\lambda\rangle\bigr)\\
    &=\bigl\|a|\lambda\rangle\bigr\|^2.
\end{align*}
Thus $\lambda\geq0$ (all eigenvalues are non-negative),
\& $\lambda=0$ iff $a|\lambda\rangle=0$.

Now consider
\begin{align*}
    N\bigl(a^\dagger|\lambda\rangle\bigr)
    &=\bigl([N,a^\dagger]+a^\dagger N\bigr)|\lambda\rangle\\
    &=\bigl(a^\dagger+a^\dagger\lambda\bigr)|\lambda\rangle\\
    &=(\lambda+1)\bigl(a^\dagger|\lambda\rangle\bigr).
\end{align*}
Similarly,
\begin{align*}
    N\bigl(a|\lambda\rangle\bigr)
    &=\bigl([N,a]+aN\bigr)|\lambda\rangle\\
    &=(-a+a\lambda)|\lambda\rangle\\
    &=(\lambda-1)a|\lambda\rangle.
\end{align*}
Here we have used the results for commutators
\[
    [N,a^\dagger]=a^\dagger
    \qquad
    \bigl(=[a^\dagger a,a^\dagger]
    =a^\dagger[a,a^\dagger]\bigr),
\]
\&
\[
    [N,a]=-a.
\]

New states will be eigenstates provided their norms are non-zero.
Compute:
\[
    \bigl\|a|\lambda\rangle\bigr\|^2=\lambda,
\]
\begin{align*}
    \bigl\|a^\dagger|\lambda\rangle\bigr\|^2
    &=\langle\lambda|aa^\dagger|\lambda\rangle
    =\langle\lambda|a^\dagger a+1|\lambda\rangle\\
    &=\lambda+1.
\end{align*}
Therefore $a^\dagger$ \& $a$ are called
\underline{creation} (or raising) \&
\underline{annihilation} (or lowering) operators, respectively.
Since all eigenvalues are non-negative, any non-integer
$\lambda$ would eventually lead to a negative eigenvalue
under repeated lowering, which is impossible.
Thus, $\lambda$ must be a non-negative integer.
Repeated lowering then reaches the ground state
$|0\rangle$, where $a|0\rangle=0$.
Since the raising operator $a^\dagger$ can act indefinitely,
all non-negative integers are eigenvalues. Hence,
\[
\boxed{\lambda=0,1,2,\ldots}
\]
\begin{center}
    \textbf{End of lecture 2}
\end{center}

From calculations of norms above, we can choose normalised eigenstates
$|n\rangle$ related by
\[
    a^\dagger|n\rangle=\sqrt{n+1}\,|n+1\rangle,
    \qquad
    a|n\rangle=\sqrt{n}\,|n-1\rangle.
\]
Starting from any state $|n\rangle$ we can generate all others.
In particular, starting from the ground state $|0\rangle$ characterised by
\[
    a|0\rangle=0,
\]
have normalised eigenstate
\[
    |n\rangle=\frac{1}{\sqrt{n!}}(a^\dagger)^n|0\rangle.
\]
$a$, $a^\dagger$ also called \underline{ladder operators}.

\medskip

\noindent
(b) \underline{Importance of the oscillator}

\begin{itemize}
    \item A smooth potential $V(x)$ with equilibrium point $x=x_0$
    (so that $V'(x_0)=0$) can be expanded
    \[
        V(x_0+x)=V(x_0)+\frac{1}{2}V''(x_0)x^2+O(x^3).
    \]
    So, if $x=x_0$ is \underline{stable}, can set $V''(x_0)=m\omega^2$
    \& for displacements $x$ not too large, expect system to be
    well-approximated by oscillator.

    e.g.\ diatomic molecules --- quantisation of vibrational energies important
    in understanding internal energy \& heat capacities --- macroscopic consequences.

    \item More complicated systems can be analysed in terms of
    \underline{normal modes}: coherent motion in which all degrees of freedom
    oscillate with same common frequency $\omega$.
    General solution for \underline{classical} oscillator is
    \[
        x(t)=Ae^{-i\omega t}+A^*e^{+i\omega t}
        \qquad A\text{ complex constant.}
    \]
    Normal modes for a system with variables $x_N$ are solutions of form
    \[
        x_N(t)=Au_Ne^{-i\omega t}+A^*u_N^*e^{+i\omega t}
    \]
    with $u_N$ complex vector specific to each normal frequency $\omega$.

    e.g.\ crystal with $\#\text{ atom }\mathcal{N}\sim10^{23}$ --- forces
    approximately elastic \& $N$ labels directions of motion for each atom,
    so $3\times\mathcal{N}$ values.

    \underline{or} EM field --- normal modes oscillations of electric \& magnetic
    fields at each point in space
    \[
        \underline{E}(\underline{x},t)
        =A\underline{u}(\underline{x})e^{-i\omega t}
        +A^*\underline{u}(\underline{x})^*e^{+i\omega t}
    \]
    for some complex vector field $\underline{u}(\underline{x})$.
    This is \underline{exact}.

    \item QM: normal modes are independent oscillators --- treat these as above.
    For each mode have
    \[
        a(\omega)\quad\&\quad a(\omega)^\dagger.
    \]
    Ground state $|0\rangle$ with
    \[
        a(\omega)|0\rangle=0.
    \]
    Other states e.g.\ $a^\dagger(\omega)|0\rangle$ \& find they have energy
    \& momentum \& behave like \underline{particles}.
    \[
        \left.
        \begin{aligned}
            a&\text{ destroys}\\
            a^\dagger&\text{ creates}
        \end{aligned}
        \right\}
        \text{particles}
        \left\{
        \begin{aligned}
            &\text{phonons in crystal}\\
            &\text{photons in EM field}
        \end{aligned}
        \right.
    \]
\end{itemize}

\subsection*{1.3 Bases \& Measurements}

\noindent
(a) \underline{Bases of Eigenstates}

From key properties (i), (ii), (iii) in \S\,1.1 we know that
for any observable $Q$ there is an orthonormal basis of eigenstates
$\{|n\rangle\}$ for $V$
\[
    Q|n\rangle=q_n|n\rangle,
    \qquad
    \langle m|n\rangle=\delta_{mn}.
\]
For any state,
\[
    |\psi\rangle=\sum_n\alpha_n|n\rangle,
    \qquad\text{where}\quad
    \alpha_n=\langle n|\psi\rangle.
\]
Note
\[
    \bigl\|\,|\psi\rangle\,\bigr\|^2=1
    \quad\Longleftrightarrow\quad
    \begin{aligned}
        \langle\psi|\psi\rangle
        &=\left(\sum_m\alpha_m^*\langle m|\right)
        \left(\sum_n\alpha_n|n\rangle\right)\\
        &=\sum_n|\alpha_n|^2=1.
    \end{aligned}
\]
Note that we can write
\[
    1=\sum_n|n\rangle\langle n|
    \qquad\text{``completeness relation''}
\]
since on any state
\begin{align*}
    |\psi\rangle
    &=\left(\sum_n|n\rangle\langle n|\right)|\psi\rangle\\
    &=\sum_n|n\rangle\bigl(\langle n|\psi\rangle\bigr),
\end{align*}
so holds for any orthonormal basis.

In above there may be several independent states with same eigenvalue $\lambda$.
The \underline{eigenspace} for a given eval is
\[
    V_\lambda=\bigl\{|\psi\rangle:Q|\psi\rangle=\lambda|\psi\rangle\bigr\}.
\]
\& note $V_\lambda$ has basis $\bigl\{|n\rangle:q_n=\lambda\bigr\}$.
The \underline{degeneracy} of $\lambda$ is the number of states in this basis,
i.e.\ $\dim V_\lambda$, but note that $\lambda$ is also called
\underline{non-degenerate} if degeneracy is $1$.

\begin{center}
    \textbf{End of lecture 3}
\end{center}

\noindent
(b) \underline{Functions of Operators}

It is straightforward to define $f(Q)$ for some operator $Q$
if $f$ is a polynomial or power series (assuming convergence), e.g.
\[
    f(Q)=Q^3+Q
    \qquad\text{or}\qquad
    f(Q)=\exp Q=\sum_{n=0}^{\infty}\frac{1}{n!}Q^n.
\]
Important examples are
\[
    e^{-iHt/\hbar}
    \qquad\text{time evolution operator ($H$ Hamiltonian)}
\]
\&
\[
    e^{-ia\hat{p}/\hbar}
    \qquad\text{translation operator (see later section of Chapter 1).}
\]

More generally, for $Q$ an observable with basis of eigenstates as in (a),
we can define
\[
    f(Q)|n\rangle=f(q_n)|n\rangle
\]
\& then defined on any state, provided $f(q_n)$ makes sense, e.g.
\[
    \frac{1}{Q}|n\rangle=\frac{1}{q_n}|n\rangle
    \qquad\text{defined for }q_n\neq0.
\]
In addition, note that
\[
    Q=\sum_n q_n|n\rangle\langle n|
    \qquad\&\qquad
    f(Q)=\sum_n f(q_n)|n\rangle\langle n|.
\]
Recall, introduced eigenspaces
\[
    V_\lambda=\bigl\{|\psi\rangle:Q|\psi\rangle=\lambda|\psi\rangle\bigr\}.
\]
A \underline{projection operator} $P_\lambda$ onto this subspace,
$V\to V_\lambda$, can be defined by its action on basis states:
\[
    P_\lambda|n\rangle=
    \begin{cases}
        |n\rangle & \text{if }q_n=\lambda,\\
        0 & \text{else}.
    \end{cases}
\]
This can be written
\[
    P_\lambda=\sum_{n:\,q_n=\lambda}|n\rangle\langle n|.
\]

\medskip
\noindent
(c) \underline{Measurement Axioms}

Consider a measurement of $Q$ (observable) when the system is in state
$|\psi\rangle$. Then
\begin{itemize}
    \item The result is an eigenvalue $\lambda$ of $Q$.
    \item This result is obtained with probability
    \[
        p(\lambda)=\bigl\|P_\lambda|\psi\rangle\bigr\|^2
        =\sum_{n:\,q_n=\lambda}|\alpha_n|^2.
    \]
    \item Immediately after the measurement the state is
    \[
        cP_\lambda|\psi\rangle
        =c\sum_{n:\,q_n=\lambda}\alpha_n|n\rangle,
    \]
    with $c$ a constant for normalisation.
\end{itemize}

Often have case with $\lambda$ non-degenerate ($\dim V_\lambda=1$),
eigenstate $|n\rangle$ unique (up to a phase), \& then
\begin{itemize}
    \item $p(\lambda)=|\alpha_n|^2=|\langle n|\psi\rangle|^2$,
    \& $\alpha_n=\langle n|\psi\rangle$ is called the \underline{amplitude}.
    \item State after measurement is $|n\rangle$.
\end{itemize}
In general
\[
    \sum_\lambda p(\lambda)=\sum_n|\alpha_n|^2=1.
\]

\medskip
\noindent
\underline{Example}

System with space $V$ with basis states (orthonormal)
\[
    |1\rangle,\quad |2\rangle,\quad |3\rangle,
\]
\& these are eigenstates of $Q$ with evals
\[
    q_1=q_2=0,\qquad q_3=1.
\]
State of system
\[
    |\psi\rangle=\frac{1}{\sqrt{6}}\bigl(2|1\rangle+|2\rangle+|3\rangle\bigr).
\]
\[
\begin{array}{c@{\qquad}c@{\qquad}c}
    \text{Result} & \text{Probability of measuring} & \text{State afterwards}\\[6pt]
    1 & \left|\dfrac{1}{\sqrt{6}}\right|^2=\dfrac{1}{6}
      & |3\rangle\\[12pt]
    0 & \left|\dfrac{2}{\sqrt{6}}\right|^2
        +\left|\dfrac{1}{\sqrt{6}}\right|^2=\dfrac{5}{6}
      & \dfrac{1}{\sqrt{5}}\bigl(2|1\rangle+|2\rangle\bigr)
\end{array}
\]

\medskip
\underline{Expectation value} (mean) of $Q$ in state $|\psi\rangle$ is
\[
    \langle Q\rangle_\psi=\langle\psi|Q|\psi\rangle
    =\sum_\lambda\lambda p(\lambda)=\sum_n q_n|\alpha_n|^2,
\]
\& the \underline{uncertainty} (spread or variance) is
\[
    (\Delta Q)_\psi^2
    =\bigl\langle(Q-\langle Q\rangle_\psi)^2\bigr\rangle_\psi
    =\langle Q^2\rangle_\psi-\langle Q\rangle_\psi^2.
\]
In thinking about many repeated measurements: must prepare the system
in the same state $|\psi\rangle$ each time.

\medskip
\noindent
Note: If $|\psi\rangle$ is an eigenstate of $Q$ with eval $\lambda$, then
\[
    \langle Q\rangle_\psi=\lambda,
    \qquad\&\qquad (\Delta Q)_\psi=0.
\]
QM behaviour arises from fact that certain observables do not commute.
In general, in any state $|\psi\rangle$,
\[
    (\Delta A)_\psi(\Delta B)_\psi
    \geq\frac{1}{2}\bigl|\langle[A,B]\rangle_\psi\bigr|,
\]
so $[A,B]\neq0$ $\Rightarrow$ cannot expect to measure $A$ \& $B$
precisely at same time.

\noindent
[Result follows from
\[
    \bigl\|(A+i\lambda B)|\psi\rangle\bigr\|^2\geq0
\]
for all real $\lambda$; expand \& look at discriminant of quadratic in $\lambda$.]

Prototype case: position \& momentum $\hat{x}$, $\hat{p}$ obey
\[
    [\hat{x},\hat{p}]=i\hbar\qquad\text{in one dimension},
    \qquad\&\qquad \Delta x\,\Delta p\geq\hbar/2.
\]
In three dimensions have $\hat{x}_i$ \& $\hat{p}_i$ with
\[
    [\hat{x}_i,\hat{p}_j]=i\hbar\delta_{ij},
\]
\& uncertainty applies to position/momentum in a given direction.

\medskip
\noindent
(d) \underline{Simultaneous measurement \& complete commuting sets}

If observable $Q$ does not distinguish completely, because some evals
are degenerate, then we can look for an additional observable $Q'$
such that $Q$, $Q'$ are simultaneously measurable.
This means we have joint eigenstates $|\lambda,\lambda'\rangle$ with
\begin{align*}
    Q|\lambda,\lambda'\rangle&=\lambda|\lambda,\lambda'\rangle,\\
    Q'|\lambda,\lambda'\rangle&=\lambda'|\lambda,\lambda'\rangle.
\end{align*}
This implies
\[
    QQ'|\lambda,\lambda'\rangle=Q'Q|\lambda,\lambda'\rangle.
\]
Hence, if we have a basis of such eigenstates, then $QQ'=Q'Q$ on any
state, or $[Q,Q']=0$.
Conversely, if $[Q,Q']=0$, easy to check that
\[
    Q':V_\lambda\to V_\lambda
    \qquad(V_\lambda\text{ eigenspace for }Q),
\]
\& finding a basis for $V_\lambda$ consisting of eigenstates of $Q'$
then gives us a basis of joint eigenstates.

\begin{center}
    \textbf{End of lecture 4}
\end{center}

\noindent
\S\,1.3 (d) --- continued

A set of observables $\{Q,Q',Q'',\ldots\}$ is a \underline{complete}
\underline{commuting} set if any pair commute $\Rightarrow$ a basis of joint
eigenstates, \& the set of eigenvalues, $\lambda,\lambda',\lambda'',\ldots$
label states \underline{uniquely} (up to normalisation).

Equivalent, if $A$ is any observable such that
\[
    [A,Q]=[A,Q']=\cdots=0
\]
then
\[
    A=f(Q,Q',\ldots).
\]

\noindent
\underline{Example}

Refer to example in (c): orthonormal states $|1\rangle$, $|2\rangle$, $|3\rangle$
eigenstates of $Q$ with $\lambda(=q_1=q_2)=0$, $\lambda(=q_3)=1$.

Second observable $Q'$ with
\[
    \begin{aligned}
        Q'|1\rangle&=|2\rangle,\\
        Q'|2\rangle&=|1\rangle,\\
        Q'|3\rangle&=0.
    \end{aligned}
\]
Check
\[
    [Q,Q']=0
    \qquad (Q':V_{\lambda=0}\to V_{\lambda=0}).
\]
Let
\[
    |\pm\rangle=\frac{1}{\sqrt{2}}\bigl(|1\rangle\pm|2\rangle\bigr)
\]
then
\[
    Q'|\pm\rangle=\pm|\pm\rangle.
\]
$Q,Q'$ complete commuting set with basis of joint eigenstates
\[
    \begin{array}{c@{\qquad}rr}
        &\lambda&\lambda'\\
        |+\rangle&0&1\\
        |-\rangle&0&-1\\
        |3\rangle&1&0
    \end{array}
\]

\subsection*{1.4 Bases, Representations \& Wavefunctions}

\noindent
(a) \underline{Discrete \& Continuous Bases}

Consider orthonormal basis $\{|n\rangle\}$ with $n$ discrete label.
We can regard all states \& operators as row/col vecs \& matrices, respectively:
\[
    \begin{aligned}
        |\psi\rangle&=\sum_n\alpha_n|n\rangle
        &\quad\Longleftrightarrow\quad \alpha_n&=\langle n|\psi\rangle,\\
        |\phi\rangle&=\sum_n\beta_n|n\rangle
        &\quad\Longleftrightarrow\quad \beta_n&=\langle n|\phi\rangle.
    \end{aligned}
    \qquad\text{cmpts of complex vectors}
\]
Inner product
\[
    \langle\phi|\psi\rangle=\sum_n\beta_n^*\alpha_n.
\]
For an operator $A$,
\[
    A|n\rangle=\sum_m|m\rangle A_{mn},
    \qquad\text{complex matrix }A_{mn}=\langle m|A|n\rangle.
\]
Simple to check that
\[
    |\phi\rangle=A|\psi\rangle
    \quad\Longleftrightarrow\quad
    \beta_m=\sum_n A_{mn}\alpha_n
\]
\& also, if $B$ is another operator, matrix $B_{mn}=\langle m|B|n\rangle$, then
\[
    (AB)_{mn}=\sum_p A_{mp}B_{pn}
    \qquad\text{(matrix multiplication)}.
\]
The choice of basis provides a representation of the quantum states \& operators.

We can also allow the index $n$ to be continuous. This requires modifications
\[
    \sum_n\ \longrightarrow\ \int dn
\]
\[
    \langle n|m\rangle=\delta_{nm}\ \longrightarrow\ \delta(n-m)
\]
\[
    |\psi\rangle=\sum_n\alpha_n|n\rangle
    \ \longrightarrow\ \int dn\,\alpha_n|n\rangle
\]
\[
    \alpha_n=\langle n|\psi\rangle.
\]
Completeness relation
\[
    1=\sum_n|n\rangle\langle n|
    \ \longrightarrow\ \int dn\,|n\rangle\langle n|.
\]
Note: if basis states above arise as eigenstates of observable, then
\[
    |\alpha_n|^2=|\langle n|\psi\rangle|^2
    \ \longrightarrow\ \text{probability \underline{density} for continuous range of }n.
\]

\medskip
\noindent
(b) \underline{Position \& Momentum Bases --- Wavefunctions}

Consider a particle in one dimension, position $\hat{x}$ \& momentum $\hat{p}$
obeying
\[
    [\hat{x},\hat{p}]=i\hbar.
\]
Let $|x\rangle$ be position eigenstates
\[
    \hat{x}|x\rangle=x|x\rangle
\]
continuous evals with
\[
    \langle x|x'\rangle=\delta(x-x'),
    \qquad \int dx\,|x\rangle\langle x|=1.
\]
Any state can be written
\[
    |\psi\rangle=\int dx\,\psi(x)|x\rangle
\]
with
\[
    \psi(x)=\langle x|\psi\rangle
\]
\& this is the usual position space wavefunction.
From remarks in (a), $|\psi(x)|^2$ is a probability density for measuring position.

Inner product
\begin{align*}
    \langle\phi|\psi\rangle
    &=\left(\int dx\,\phi(x)^*\langle x|\right)
      \left(\int dx'\,\psi(x')|x'\rangle\right)\\
    &=\int dx\int dx'\,\phi(x)^*\psi(x')
      \underbrace{\langle x|x'\rangle}_{\delta(x-x')}
      =\int dx\,\phi(x)^*\psi(x)
\end{align*}
\&
\[
    \bigl\|\,|\psi\rangle\,\bigr\|^2
    =\langle\psi|\psi\rangle=\int dx\,|\psi(x)|^2=1,
    \quad\text{if normalised}.
\]
Similarly, define momentum eigenstates $|p\rangle$ with
\[
    \hat{p}|p\rangle=p|p\rangle
\]
\&
\[
    \langle p|p'\rangle=\delta(p-p'),
    \qquad \int dp\,|p\rangle\langle p|=1.
\]
Now can expand
\[
    |\psi\rangle=\int dp\,\tilde{\psi}(p)|p\rangle
\]
with
\[
    \tilde{\psi}(p)=\langle p|\psi\rangle
\]
momentum space wavefunction \& $|\tilde{\psi}(p)|^2$ probability distribution
for measurements of momentum.

To proceed, we need
\[
    \left.
    \begin{aligned}
        \langle x|p\rangle&=\frac{1}{(2\pi\hbar)^{1/2}}e^{ipx/\hbar},\\
        \langle p|x\rangle&=\frac{1}{(2\pi\hbar)^{1/2}}e^{-ipx/\hbar}.
    \end{aligned}
    \right\}\qquad (*)
\]
Assume $(*)$ holds \& verify it below.

Consider action of $\hat{x}$ \& $\hat{p}$ on position space wavefunctions
\[
    \hat{x}|\psi\rangle\quad\text{wavefn}:\quad
    \langle x|\hat{x}|\psi\rangle=x\langle x|\psi\rangle=x\psi(x).
\]
\begin{align*}
    \hat{p}|\psi\rangle\quad\text{wavefn}:\quad
    &\langle x|\hat{p}|\psi\rangle\\
    &=\int dp\,\langle x|\hat{p}|p\rangle\langle p|\psi\rangle
    \qquad\text{completeness}\\
    &=\int dp\,p\langle x|p\rangle\langle p|\psi\rangle.
\end{align*}
Note
\[
    p\langle x|p\rangle=-i\hbar\frac{\partial}{\partial x}\langle x|p\rangle
    \qquad\text{from }(*).
\]
Hence
\begin{align*}
    \langle x|\hat{p}|\psi\rangle
    &=-i\hbar\frac{\partial}{\partial x}
      \int dp\,\langle x|p\rangle\langle p|\psi\rangle\\
    &=-i\hbar\frac{\partial}{\partial x}\langle x|\psi\rangle
      \qquad\text{completeness again}\\
    &=-i\hbar\psi'(x).
\end{align*}
i.e.
\[
    \left.
    \begin{aligned}
        \hat{x}&\to x\\
        \hat{p}&\to-i\hbar\,\partial/\partial x
    \end{aligned}
    \right\}\quad\text{on position space wavefunctions}.
\]
Similarly, on momentum space wavefunctions, calculate
\[
    \begin{gathered}
        \langle p|\hat{x}|\psi\rangle\\
        \langle p|\hat{p}|\psi\rangle
    \end{gathered}
    \qquad\text{\& find}\qquad
    \begin{aligned}
        \hat{x}&\to i\hbar\frac{\partial}{\partial p}\\
        \hat{p}&\to p.
    \end{aligned}
\]
Relationship between wavefunctions $\psi(x)$ \& $\tilde{\psi}(p)$ also follows from $(*)$
\begin{align*}
    \tilde{\psi}(p)&=\langle p|\psi\rangle
    =\int dx\,\langle p|x\rangle\langle x|\psi\rangle\\
    &=\frac{1}{(2\pi\hbar)^{1/2}}\int dx\,e^{-ipx/\hbar}\psi(x)
\end{align*}
\&
\begin{align*}
    \psi(x)&=\langle x|\psi\rangle
    =\int dp\,\langle x|p\rangle\langle p|\psi\rangle\\
    &=\frac{1}{(2\pi\hbar)^{1/2}}\int dp\,e^{ixp/\hbar}\tilde{\psi}(p).
\end{align*}
Fourier transform \& inverse.

\begin{center}
    \textbf{End of lecture 5}
\end{center}

Now given $H(\hat{x},\hat{p})$ we can consider
\[
    H|\psi\rangle=E|\psi\rangle.
\]
If
\[
    H=\frac{\hat{p}^2}{2m}+V(\hat{x})
\]
\[
    \begin{aligned}
        \text{on }\psi(x)\qquad H&\to-\frac{\hbar^2}{2m}\frac{\partial^2}{\partial x^2}+V(x),\\
        \text{on }\tilde{\psi}(p)\qquad H&\to\frac{p^2}{2m}+V\!\left(i\hbar\frac{\partial}{\partial p}\right).
    \end{aligned}
\]
It may be easy to interpret potential term in momentum space.
e.g.\ $V(x)=\lambda x^n$ then
\[
    V\!\left(i\hbar\frac{\partial}{\partial p}\right)
    =\lambda(i\hbar)^n\left(\frac{\partial}{\partial p}\right)^n
\]
for $n$ a positive integer.
More generally, argue from first principles as follows.

Consider
\begin{align*}
    \langle p|V(\hat{x})|\psi\rangle
    &=\int dx\,\langle p|V(\hat{x})|x\rangle\langle x|\psi\rangle\\
    &=\int dx\int dp'\,\langle p|V(\hat{x})|x\rangle
      \langle x|p'\rangle\langle p'|\psi\rangle\\
    &=\int dp'\left(\int dx\,V(x)
      \underbrace{\langle p|x\rangle\langle x|p'\rangle}_{\text{use }(*)}
      \langle p'|\psi\rangle\right)\\
    &=\int dp'\int dx\,V(x)\frac{e^{-i(p-p')x/\hbar}}{2\pi\hbar}\tilde{\psi}(p')\\
    &=\frac{1}{(2\pi\hbar)^{1/2}}\int dp'\,\tilde{V}(p-p')\tilde{\psi}(p'),
\end{align*}
\hfill reproducing convolution theorem.

i.e.\ $H|\psi\rangle=E|\psi\rangle$ becomes
\[
    -\frac{\hbar^2}{2m}\psi''(x)+V(x)\psi(x)=E\psi
    \qquad\text{in position space}
\]
or
\[
    \frac{p^2}{2m}\tilde{\psi}(p)
    +\frac{1}{(2\pi\hbar)^{1/2}}\int dp'\,\tilde{V}(p-p')\tilde{\psi}(p')
    =E\tilde{\psi}(p)
\]
\hfill in momentum space.

To generalise from one to three dimensions: position \& momentum operators obey
\[
    [\hat{x}_i,\hat{p}_j]=i\hbar\delta_{ij}.
\]
One complete commuting set is
\[
    \underline{\hat{x}}=(\hat{x}_1,\hat{x}_2,\hat{x}_3)
\]
with joint eigenstates
\[
    \hat{x}_i|\underline{x}\rangle=x_i|\underline{x}\rangle
\]
orthonormal
\[
    \langle\underline{x}|\underline{x}'\rangle
    =\delta^{(3)}(\underline{x}-\underline{x}').
\]
An alternative set
\[
    \underline{\hat{p}}=(\hat{p}_1,\hat{p}_2,\hat{p}_3)
\]
with joint eigenstates
\[
    \hat{p}_i|\underline{p}\rangle=p_i|\underline{p}\rangle
    \qquad\text{(also orthonormal)}.
\]
For any state $|\psi\rangle$ we can expand
\[
    |\psi\rangle=\int d^3\underline{x}\,\psi(\underline{x})|\underline{x}\rangle
    =\int d^3\underline{p}\,\tilde{\psi}(\underline{p})|\underline{p}\rangle
\]
\&
\[
    \psi(\underline{x})=\langle\underline{x}|\psi\rangle,
    \qquad\tilde{\psi}(\underline{p})=\langle\underline{p}|\psi\rangle.
\]
$(*)$ becomes
\[
    \langle\underline{x}|\underline{p}\rangle
    =\frac{1}{(2\pi\hbar)^{3/2}}e^{i\underline{p}\cdot\underline{x}/\hbar}
\]
\& wavefunctions related by
\[
    \tilde{\psi}(\underline{p})=\frac{1}{(2\pi\hbar)^{3/2}}
    \int d^3\underline{x}\,e^{-i\underline{p}\cdot\underline{x}/\hbar}\psi(\underline{x})
    \qquad\text{(\& inverse)}.
\]

\medskip
\noindent
(c) \underline{More on position \& momentum eigenstates}

Return to one dimension \& seek to justify $(*)$ using only
$[\hat{x},\hat{p}]=i\hbar$.

Given $|x\rangle$ with $\hat{x}|x\rangle=x|x\rangle$, consider
\[
    T(a)=e^{-ia\hat{p}/\hbar}
    =\sum_{n=0}^{\infty}\frac{1}{n!}\left(\frac{-ia}{\hbar}\right)^n\hat{p}^n
    \qquad a\text{ any real number}
\]
\& note
\[
    [\hat{x},T(a)]=\left[\hat{x},\frac{-ia\hat{p}}{\hbar}\right]T(a)=aT(a)
    \qquad\text{[see 3 on Example Sheet 1]}.
\]
Hence
\begin{align*}
    \hat{x}\bigl(T(a)|x\rangle\bigr)
    &=\bigl([\hat{x},T(a)]+T(a)\hat{x}\bigr)|x\rangle\\
    &=(x+a)\bigl(T(a)|x\rangle\bigr).
\end{align*}
Think of $T(a)|x\rangle$ as a translated eigenstate.

Similarly, consider
\[
    S(b)=e^{ib\hat{x}/\hbar}
\]
\& note
\[
    [\hat{p},S(b)]=bS(b).
\]
Hence
\[
    \hat{p}|p\rangle=p|p\rangle
    \quad\Rightarrow\quad
    \hat{p}\bigl(S(b)|p\rangle\bigr)=(p+b)\bigl(S(b)|p\rangle\bigr).
\]
$S(b)$ gives translation in momentum space.

Now take fixed values $x_0$ \& $p_0$ \& define all eigenstates by
\[
    |x_0+a\rangle=T(a)|x_0\rangle,
    \qquad |p_0+b\rangle=S(b)|p_0\rangle.
\]
Now observe
\begin{align*}
    \langle x_0+a|p_0+b\rangle
    &=\langle x_0|e^{ia\hat{p}/\hbar}|p_0+b\rangle\\
    &=\langle x_0|e^{ia(p_0+b)/\hbar}e^{ib\hat{x}/\hbar}|p_0\rangle\\
    &=e^{i(ab+ap_0+bx_0)/\hbar}
      \underbrace{\langle x_0|p_0\rangle}_{\displaystyle\text{choose }=\frac{1}{\sqrt{2\pi\hbar}}}.
\end{align*}
\hfill for correct, consistent normalisation.

Then set $x_0=p_0=0$ to get
\[
    \langle a|b\rangle=\frac{1}{\sqrt{2\pi\hbar}}e^{iab/\hbar}
    \qquad\text{which is }(*).
\]

\section*{2. Evolution in Time \& Pictures}
\subsection*{2.1 The Schr\"odinger Equation \& the Schr\"odinger Picture}

The evolution in time of a state $|\psi(t)\rangle$ is governed
(in absence of external influences or measurements) by the Schr\"odinger equation (SE)
\[
    i\hbar\frac{\partial}{\partial t}|\psi(t)\rangle=H|\psi(t)\rangle
\]
where $H=H^\dagger$ is the Hamiltonian.

Equivalently
\[
    -i\hbar\frac{\partial}{\partial t}\langle\psi(t)|=\langle\psi(t)|H.
\]
This implies
\[
    \frac{\partial}{\partial t}\bigl(\langle\psi(t)|\psi(t)\rangle\bigr)=0
\]
so that normalisation of $|\psi(t)\rangle$ is preserved in time.

Solution of SE can be written
\[
    |\psi(t)\rangle=U(t)|\psi(0)\rangle
\]
where
\[
    U(t)=e^{-iHt/\hbar}
    \qquad\text{time evolution operator, assuming }H\text{ independent of }t.
\]
Check immediately from
\[
    i\hbar\frac{\partial U(t)}{\partial t}=HU(t).
\]
Recall that time-evolution problems can also be conveniently solved by
working with basis of energy eigenstates $\{|n\rangle\}$
with
\[
    H|n\rangle=E_n|n\rangle
\]
with
\[
    e^{-iE_nt/\hbar}|n\rangle
    \qquad\text{solution of (t-dep) SE}.
\]

\begin{center}
    \textbf{End of lecture 6}
\end{center}

\noindent
\S\,2.1 Sch. eqn \& Sch picture

\[
    i\hbar\frac{\partial}{\partial t}|\psi(t)\rangle=H|\psi\rangle
\]
\[
    |\psi(t)\rangle=U(t)|\psi(0)\rangle
\]
where
\begin{align*}
    U(t)&=e^{-iHt/\hbar}
    \qquad(H\text{ indep of }t)\\
    &=\sum_{n=0}^{\infty}\frac{1}{n!}
      \left(-\frac{it}{\hbar}H\right)^n
\end{align*}

If
\[
    |\psi(0)\rangle=\sum_n\alpha_n|n\rangle
\]
where
\[
    H|n\rangle=E_n|n\rangle
\]
then
\[
    |\psi(t)\rangle=\sum_n\alpha_n e^{-iE_nt/\hbar}|n\rangle
\]

\medskip
\noindent
\underline{Example}\qquad system with two energy eigenstates
\[
    \begin{array}{cc}
        |1\rangle&|2\rangle\\
        E_1&E_2
    \end{array}
\]
Consider measuring $Q$ defined by
\[
    Q|1\rangle=|2\rangle,\qquad Q|2\rangle=|1\rangle
\]
So
\[
    [Q,H]\neq0
\]
Eigenstates of $Q$:
\[
    |\pm\rangle=\frac{1}{\sqrt{2}}\bigl(|1\rangle\pm|2\rangle\bigr)
    \qquad\text{evals }q_\pm=\pm1
\]
Take initial state $|\psi(0)\rangle=|+\rangle$
\[
    \Rightarrow\qquad
    |\psi(t)\rangle=\frac{1}{\sqrt{2}}
    \left(e^{-iE_1t/\hbar}|1\rangle+e^{-iE_2t/\hbar}|2\rangle\right)
\]
Probability of measuring $Q$ at time $t$ \& getting $\pm1$ is
\begin{align*}
    &|\langle\pm|\psi(t)\rangle|^2\\
    &\quad=\left|\frac{1}{2}(\langle1|\pm\langle2|)
    \left(e^{-iE_1t/\hbar}|1\rangle+e^{-iE_2t/\hbar}|2\rangle\right)\right|^2\\
    &\quad=\left|\frac{1}{2}\left(e^{-iE_1t/\hbar}\pm e^{-iE_2t/\hbar}\right)\right|^2
    =\left\{\begin{aligned}
        &\cos^2\bigl((E_1-E_2)t/2\hbar\bigr)\\
        &\sin^2\bigl((E_1-E_2)t/2\hbar\bigr)
    \end{aligned}\right.
\end{align*}

In basis
\[
    \begin{aligned}
        |1\rangle&\to\begin{pmatrix}1\\0\end{pmatrix}\\
        |2\rangle&\to\begin{pmatrix}0\\1\end{pmatrix}
    \end{aligned}
    \qquad H\to\begin{pmatrix}E_1&0\\0&E_2\end{pmatrix}
\]
\[
    U=e^{-iHt/\hbar}\quad\to\quad
    \begin{pmatrix}e^{-iE_1t/\hbar}&0\\0&e^{-iE_2t/\hbar}\end{pmatrix}
\]
In alternative basis
\[
    \begin{aligned}
        |+\rangle&\to\begin{pmatrix}1\\0\end{pmatrix}\\
        |-\rangle&\to\begin{pmatrix}0\\1\end{pmatrix}
    \end{aligned}
\]
\[
    H|\pm\rangle=\frac{1}{\sqrt{2}}\bigl(E_1|1\rangle+E_2|2\rangle\bigr)
\]
\[
    \left.\begin{aligned}
        \langle\pm|H|\pm\rangle&=\frac{1}{2}(E_1+E_2)=\mathcal{E}\\
        \langle\mp|H|\pm\rangle&=\frac{1}{2}(E_1-E_2)=\hbar\lambda
    \end{aligned}\right\}\quad\text{def }\mathcal{E}\ \&\ \lambda\text{ consts.}
\]
\[
    H\to\begin{pmatrix}\mathcal{E}&\hbar\lambda\\\hbar\lambda&\mathcal{E}\end{pmatrix}
    =\mathcal{E}I+\hbar\lambda M
\]
\[
    I\quad 2\times2\text{ identity}
    \qquad\&\qquad M=\begin{pmatrix}0&1\\1&0\end{pmatrix}
\]
Now
\begin{align*}
    U(t)=\exp(-iHt/\hbar)&\to e^{-\frac{i}{\hbar}t(\mathcal{E}I+\lambda\hbar M)}\\
    &=e^{-\frac{i}{\hbar}t\mathcal{E}}e^{-it\lambda M}
\end{align*}
Note $M^2=I$ \& hence
\begin{align*}
    \exp(-it\lambda M)
    &=\sum_{r=0}^{\infty}\frac{1}{(2r)!}(-i\lambda t)^{2r}I
      +\sum_{r=0}^{\infty}\frac{1}{(2r+1)!}(-i\lambda t)^{2r+1}M\\
    &=(\cos\lambda t)I-i(\sin\lambda t)M.
\end{align*}
\[
    |\psi(0)\rangle=|+\rangle\to\begin{pmatrix}1\\0\end{pmatrix}
\]
\[
    |\psi(t)\rangle=e^{-iHt/\hbar}|+\rangle
    \to\left(e^{-it\mathcal{E}/\hbar}\right)
    \begin{pmatrix}\cos\lambda t\\-i\sin\lambda t\end{pmatrix}
\]
\[
    e^{-it\mathcal{E}/\hbar}
    \bigl(\cos\lambda t\,|+\rangle-i\sin\lambda t\,|-\rangle\bigr)
    \qquad\longleftarrow
\]
Now compare with previous results for prob.\ of measuring $Q$ \& finding $\pm1$

\subsection*{2.2 Heisenberg Picture \& Heisenberg Equations}

\noindent
(a) \underline{Unitary operators}

Physical predictions in QM are given by probability amplitudes,
eigenvalues, expectation values, etc., hence by expressions like
\[
    \langle\phi|\psi\rangle,\quad\langle\phi|A|\psi\rangle,\quad\text{etc.}
\]
An operator $U$ is called \underline{unitary} if
\[
    UU^\dagger=U^\dagger U=1\quad\text{or}\quad U^\dagger=U^{-1}.
\]
Given any such $U$, consider a map on states \& operators
\[
    \begin{aligned}
        |\psi\rangle&\longrightarrow |\psi'\rangle=U|\psi\rangle\\
        \langle\psi|&\longrightarrow\langle\psi'|=\langle\psi|U^\dagger\\
        A&\longrightarrow A'=UAU^\dagger
    \end{aligned}
\]
Under this, all physical predictions remain unchanged
\begin{align*}
    \text{e.g.}\qquad
    \langle\phi|\psi\rangle&\to\langle\phi'|\psi'\rangle
    =\langle\phi|U^\dagger U|\psi\rangle=\langle\phi|\psi\rangle\\
    \langle\phi|A|\psi\rangle&\to\langle\phi'|A'|\psi'\rangle
    =\langle\phi|U^\dagger UAU^\dagger U|\psi\rangle\\
    &=\langle\phi|A|\psi\rangle.
\end{align*}
Furthermore
\[
    \begin{aligned}
        C=AB&\longrightarrow C'=A'B'\\
        C=[A,B]&\longrightarrow C'=[A',B']
    \end{aligned}
\]
\[
    Q^\dagger=Q\quad\Rightarrow\quad Q'=UQU^\dagger
    \quad\text{also satisfies }(Q')^\dagger=Q'
\]
\[
    Q|\psi\rangle=\lambda|\psi\rangle
    \quad\Rightarrow\quad
    \underset{\text{same eigenvalue.}}{Q'|\psi'\rangle=\lambda|\psi'\rangle}
\]

\medskip
\noindent
(b) \underline{Changing Pictures}

In the \underline{Schr\"odinger picture}
\[
    i\hbar\frac{\partial}{\partial t}|\psi(t)\rangle=H|\psi(t)\rangle
    \qquad\left(\begin{gathered}\text{all for }H\text{ indep of }t\\\&\ H=H^\dagger\end{gathered}\right)
\]
\[
    \Rightarrow\qquad |\psi(t)\rangle=U(t)|\psi(0)\rangle
\]
where $U(t)=\exp(-iHt/\hbar)$ is unitary. In fact
\[
    U(t)^\dagger=U(t)^{-1}=U(-t)
\]
\[
    \&\qquad U(t_1)U(t_2)=U(t_1+t_2)
\]
The states depend on $t$ but observables do not, in S-picture.

To pass to Heisenberg picture we shift $t$-dependence from states to
operators as follows
\[
    \begin{array}{c@{\quad}c@{\quad}l}
        &\text{Schr\"odinger}&\text{Heisenberg}\\[6pt]
        \text{states}&|\psi(t)\rangle_S
        &|\psi\rangle_H=U^\dagger|\psi(t)\rangle_S
        =e^{itH/\hbar}|\psi(t)\rangle_S=|\psi(0)\rangle_S\\[8pt]
        \text{ops}&A_S&A_H(t)=U^\dagger A_SU=e^{itH/\hbar}A_Se^{-itH/\hbar}
    \end{array}
\]
Because the transformation $U(t)$ is unitary
\[
    {}_S\langle\phi(t)|A_S|\psi(t)\rangle_S
    ={}_H\langle\phi|A_H(t)|\psi\rangle_H
\]
Note
\[
    H_H=H_S=H
\]
To specify dynamics in H-picture, consider
\begin{align*}
    \frac{d}{dt}A_H(t)
    &=\frac{d}{dt}\left(e^{iHt/\hbar}A_Se^{-iHt/\hbar}\right)\\
    &=\frac{iH}{\hbar}e^{iHt/\hbar}A_Se^{-iHt/\hbar}
      -e^{iHt/\hbar}A_Se^{-iHt/\hbar}\frac{iH}{\hbar}\\
    &=\frac{i}{\hbar}[H,A_H(t)]
\end{align*}
\underline{or}
\[
    i\hbar\frac{d}{dt}A_H(t)=[A_H(t),H]
    \qquad\underline{\text{Heisenberg equations of motion}}
\]
e.g.\ particle in one dimension $\hat{x}(t)$, $\hat{p}(t)$ in H-picture

consider
\[
    H=\frac{\hat{p}^2}{2m}+V(\hat{x})
\]
\[
    \frac{d}{dt}\hat{x}(t)=\frac{1}{i\hbar}[\hat{x}(t),H]
    =\frac{1}{m}\hat{p}(t)
\]
\[
    \frac{d}{dt}\hat{p}(t)=\frac{1}{i\hbar}[\hat{p}(t),H]
    =-V'(\hat{x}(t))
\]
using
\[
    [\hat{x}(t),\hat{p}(t)]=i\hbar
\]

\begin{center}
    \textbf{End of lecture 7}
\end{center}

Taking expectation values in any state $|\psi\rangle$ (independent of $t$
in the Heisenberg picture) gives
\[
    \frac{d}{dt}\langle\hat{x}\rangle=\frac{1}{m}\langle\hat{p}\rangle,
    \qquad
    \frac{d}{dt}\langle\hat{p}\rangle=-\langle V'(\hat{x})\rangle.
\]
These are \underline{Ehrenfest's theorem}, true in all pictures.

For some potentials we can solve the Heisenberg equations explicitly.

\medskip
\noindent
\underline{Free particle: $V=0$}
\[
    \frac{d}{dt}\hat{p}(t)=0
    \quad\Rightarrow\quad
    \hat{p}(t)=\hat{p}(0)
    \qquad\text{(constant operator).}
\]
Then
\[
    \frac{d}{dt}\hat{x}(t)=\frac{1}{m}\hat{p}(0)
    \quad\Rightarrow\quad
    \hat{x}(t)=\hat{x}(0)+\frac{\hat{p}(0)}{m}t,
\]
just as in classical dynamics.

\medskip
\noindent
\underline{Harmonic oscillator: $V=\frac{1}{2}m\omega^2x^2$}
\[
    \frac{d}{dt}\hat{x}(t)=\frac{1}{m}\hat{p}(t),
    \qquad
    \frac{d}{dt}\hat{p}(t)=-m\omega^2\hat{x}(t).
\]
Hence
\[
    \frac{d^2}{dt^2}\hat{x}(t)+\omega^2\hat{x}(t)=0,
    \qquad
    \frac{d^2}{dt^2}\hat{p}(t)+\omega^2\hat{p}(t)=0.
\]
The solutions are
\begin{align*}
    \hat{x}(t)&=\hat{x}(0)\cos\omega t
        +\frac{\hat{p}(0)}{m\omega}\sin\omega t,\\
    \hat{p}(t)&=\hat{p}(0)\cos\omega t
        -m\omega\hat{x}(0)\sin\omega t.
\end{align*}
Or equivalently,
\begin{align*}
    \hat{x}(t)&=\left(\frac{\hbar}{2m\omega}\right)^{1/2}
        \left(ae^{-i\omega t}+a^\dagger e^{i\omega t}\right),\\
    \hat{p}(t)&=\left(\frac{\hbar m\omega}{2}\right)^{1/2}\frac{1}{i}
        \left(ae^{-i\omega t}-a^\dagger e^{i\omega t}\right),
\end{align*}
where $a$, $a^\dagger$ are defined from the Schr\"odinger-picture operators
$\hat{x}(0)$, $\hat{p}(0)$.

\medskip
\noindent
\underline{Why do the ladder operators acquire these phases?}

For the harmonic oscillator,
\[
    H=\hbar\omega\left(a^\dagger a+\frac{1}{2}\right),
    \qquad [a,a^\dagger]=1.
\]
Using $[AB,C]=A[B,C]+[A,C]B$, we find
\begin{align*}
    [a^\dagger a,a]&=a^\dagger[a,a]+[a^\dagger,a]a=-a,\\
    [a^\dagger a,a^\dagger]
        &=a^\dagger[a,a^\dagger]+[a^\dagger,a^\dagger]a=a^\dagger.
\end{align*}
The constant term in $H$ commutes with every operator, so
\[
    [H,a]=-\hbar\omega a,
    \qquad [H,a^\dagger]=\hbar\omega a^\dagger.
\]
Define
\[
    a_H(t)=U^\dagger(t)aU(t),
    \qquad a_H^\dagger(t)=U^\dagger(t)a^\dagger U(t).
\]
Since $H$ commutes with $U(t)$, the same commutator relations hold for
the Heisenberg-picture operators. Their equations of motion are therefore
\begin{align*}
    \frac{d}{dt}a_H(t)&=\frac{i}{\hbar}[H,a_H(t)]
        =-i\omega a_H(t),\\
    \frac{d}{dt}a_H^\dagger(t)&=\frac{i}{\hbar}[H,a_H^\dagger(t)]
        =i\omega a_H^\dagger(t).
\end{align*}
With $a_H(0)=a$ and $a_H^\dagger(0)=a^\dagger$, this gives
\[
    \boxed{a_H(t)=ae^{-i\omega t},\qquad
           a_H^\dagger(t)=a^\dagger e^{i\omega t}.}
\]
The two phase factors have opposite signs because the commutators of $H$
with the lowering and raising operators have opposite signs. The second
result also follows by taking the adjoint of the first.
Substituting these results into
\[
    \hat{x}(t)=\sqrt{\frac{\hbar}{2m\omega}}
        \bigl(a_H(t)+a_H^\dagger(t)\bigr),
    \qquad
    \hat{p}(t)=\frac{1}{i}\sqrt{\frac{\hbar m\omega}{2}}
        \bigl(a_H(t)-a_H^\dagger(t)\bigr)
\]
recovers the expressions above. Here $a$ and $a^\dagger$ without a time
argument denote the fixed operators at $t=0$.

\subsection*{2.3 Canonical Quantisation}

\noindent
$\ast$ \underline{Non-examinable material}

This is the final step in Dirac's approach to constructing QM:
how does one pass from a general classical system to its quantum version?

A classical system can be described by a set of generalised positions
$x_i(t)$ and momenta $p_i(t)$, $1\leq i\leq N$ (which may include angles
and angular momenta). These are coordinates on \underline{phase space}.
The dynamics are governed by a Hamiltonian $H(x_i,p_i)$.

The \underline{Poisson bracket} of two functions $f(x_i,p_i)$,
$g(x_i,p_i)$ is defined by
\[
    \{f,g\}=\sum_{i=1}^{N}\left(
        \frac{\partial f}{\partial x_i}\frac{\partial g}{\partial p_i}
        -\frac{\partial f}{\partial p_i}\frac{\partial g}{\partial x_i}
    \right),
\]
which is another function on phase space. In particular,
\[
    \{x_i,p_j\}=\delta_{ij}.
\]
Properties of Poisson brackets closely resemble those of commutators
(linearity, the Jacobi identity, etc.). Dynamics in this formulation are
given by Hamilton's equation
\[
    \frac{df}{dt}=\{f,H\}
    \qquad\text{for }f\text{ with no explicit time dependence.}
\]
In \underline{canonical quantisation}, classical functions are promoted
to quantum operators,
\[
    f,g\quad\longrightarrow\quad\hat{f},\hat{g},
\]
and Poisson brackets determine commutators through Dirac's formal rule
\[
    [\hat{f},\hat{g}]=i\hbar\,\widehat{\{f,g\}}.
\]
In particular, this recovers
\[
    [\hat{x}_i,\hat{p}_j]=i\hbar\delta_{ij},
\]
and the correspondence with Hamilton's equation gives the Heisenberg
equation
\[
    i\hbar\frac{d\hat{f}}{dt}=[\hat{f},H].
\]
The rule is a quantisation prescription; for general functions, operator
ordering must be specified and the rule cannot hold exactly for every
pair of functions simultaneously.

\noindent
$\ast$ \underline{End of non-examinable material}

\subsection*{Tensor Products}

Given state spaces $V_1$ and $V_2$, their \underline{tensor product}
$V=V_1\otimes V_2$ is spanned by tensor product states
\[
    |\psi\rangle\otimes|\phi\rangle,
    \qquad |\psi\rangle\in V_1,\quad |\phi\rangle\in V_2,
\]
with duals $\langle\psi|\otimes\langle\phi|$, subject to
\begin{align*}
    (|\psi\rangle+|\psi'\rangle)\otimes|\phi\rangle
        &=|\psi\rangle\otimes|\phi\rangle
          +|\psi'\rangle\otimes|\phi\rangle,\\
    |\psi\rangle\otimes(|\phi\rangle+|\phi'\rangle)
        &=|\psi\rangle\otimes|\phi\rangle
          +|\psi\rangle\otimes|\phi'\rangle,\\
    (\alpha|\psi\rangle)\otimes|\phi\rangle
        &=|\psi\rangle\otimes(\alpha|\phi\rangle)
          =\alpha\bigl(|\psi\rangle\otimes|\phi\rangle\bigr),
\end{align*}
and similarly for duals.

The inner product on $V$ satisfies
\[
    \bigl(\langle\psi'|\otimes\langle\phi'|\bigr)
    \bigl(|\psi\rangle\otimes|\phi\rangle\bigr)
    =\langle\psi'|\psi\rangle\langle\phi'|\phi\rangle,
\]
extended sesquilinearly to arbitrary states (conjugate linear in the
first state and linear in the second).

Given operators $A$ on $V_1$ and $B$ on $V_2$, define an operator
$A\otimes B$ on $V$ by
\[
    (A\otimes B)\bigl(|\psi\rangle\otimes|\phi\rangle\bigr)
    =(A|\psi\rangle)\otimes(B|\phi\rangle),
\]
and extend by linearity. In particular,
\[
    A\longleftrightarrow A\otimes 1
    \quad\text{acting only on }V_1,
    \qquad
    B\longleftrightarrow 1\otimes B
    \quad\text{acting only on }V_2.
\]
Operators acting on different factors commute for any $A,B$:
\[
    [A\otimes 1,1\otimes B]=0.
\]

\medskip
\noindent
\underline{Example: a particle in two dimensions}

Consider the position operators $\hat{x}_1$, $\hat{x}_2$.
A basis of joint eigenstates is constructed as
\[
    \underbrace{|x_1,x_2\rangle}_{\in V}
    =\underbrace{|x_1\rangle}_{\in V_1}
       \otimes\underbrace{|x_2\rangle}_{\in V_2},
\]
where $V_1$, $V_2$ are the state spaces for a particle in one dimension.
The position operators are identified with
\[
    \hat{x}_1\longleftrightarrow\hat{x}_1\otimes 1,
    \qquad
    \hat{x}_2\longleftrightarrow 1\otimes\hat{x}_2.
\]
The wavefunction for a product state $|\psi\rangle\otimes|\phi\rangle$
is given by
\begin{align*}
    \langle x_1,x_2|\bigl(|\psi\rangle\otimes|\phi\rangle\bigr)
        &=(\langle x_1|\otimes\langle x_2|)
          (|\psi\rangle\otimes|\phi\rangle)\\
        &=\langle x_1|\psi\rangle\langle x_2|\phi\rangle\\
        &=\psi(x_1)\phi(x_2).
\end{align*}
Thus
\[
    \begin{aligned}
        \text{tensor product states}
            &\longleftrightarrow\text{separable wavefunctions},\\
        \text{tensor product space}
            &\longleftrightarrow\text{all linear combinations of such wavefunctions}.
    \end{aligned}
\]

\medskip
\noindent
\underline{Note.}
If $\{|n\rangle\}$ is a basis for $V_1$ and $\{|m\rangle\}$ is a basis
for $V_2$, then
\[
    \{|n\rangle\otimes|m\rangle\}
\]
is a basis for $V$. If both bases are finite,
\[
    \dim V=(\dim V_1)(\dim V_2).
\]

\begin{center}
    \textbf{End of lecture 8}
\end{center}
\medskip
\noindent
\underline{Example: the two-dimensional isotropic oscillator}

The Hamiltonian is
\[
    H=\frac{1}{2m}\left(\hat{p}_1^2+\hat{p}_2^2\right)
      +\frac{1}{2}m\omega^2\left(\hat{x}_1^2+\hat{x}_2^2\right)
      =H_1+H_2,
\]
where
\[
    H_i=\frac{\hat{p}_i^2}{2m}+\frac{1}{2}m\omega^2\hat{x}_i^2
       =\hbar\omega\left(N_i+\frac{1}{2}\right),
    \qquad i=1,2.
\]
Here $N_i=a_i^\dagger a_i$, with
\[
    [a_i,a_j^\dagger]=\delta_{ij}.
\]
Since $[N_1,N_2]=0$, we construct simultaneous eigenstates
\[
    |n_1,n_2\rangle=|n_1\rangle\otimes|n_2\rangle,
    \qquad n_1,n_2=0,1,2,\ldots.
\]
Then
\begin{align*}
    H|n_1,n_2\rangle
        &=(H_1|n_1\rangle)\otimes|n_2\rangle
          +|n_1\rangle\otimes(H_2|n_2\rangle)\\
        &=E_{n_1n_2}|n_1,n_2\rangle,
\end{align*}
where
\[
    E_{n_1n_2}=\hbar\omega(n_1+n_2+1).
\]
The corresponding wavefunctions factorise:
\begin{align*}
    \psi_{n_1n_2}(x_1,x_2)
        &=\langle x_1,x_2|n_1,n_2\rangle\\
        &=\psi_{n_1}(x_1)\psi_{n_2}(x_2),
\end{align*}
where $\psi_n(x)=\langle x|n\rangle$ is the wavefunction of the
one-dimensional oscillator.

\noindent
\underline{Notation.}
It is common to suppress the tensor product symbol and write
$|n_1\rangle|n_2\rangle$ for $|n_1\rangle\otimes|n_2\rangle$.

\subsection*{3.2 Spin}

Experiment shows that particles generally carry an internal degree of
freedom called \underline{spin}, or \underline{intrinsic angular momentum}.
Even if a particle appears to be elementary or point-like, its state
space has the form
\[
    V=V_{\mathrm{space}}\otimes V_{\mathrm{spin}},
\]
with basis states
\[
    |\mathbf{x},r\rangle=|\mathbf{x}\rangle\otimes|r\rangle,
\]
where $r$ takes a finite set of values.

The position operators alone do not, in general, form a complete
commuting set: there are other observables $Q$ acting only on
$V_{\mathrm{spin}}$, with
\[
    [\hat{x}_i,Q]=0.
\]
We return to these operators in Chapter 5; for now we focus on states.

Each kind of particle has a definite total spin $s$, a non-negative
integer or half-integer,
\[
    s=0,\frac{1}{2},1,\frac{3}{2},\ldots.
\]
This is part of its nature, like mass or charge. For a massive particle
of spin $s$, the spin space has $2s+1$ basis states, labelled by
\[
    r=s,s-1,\ldots,-s.
\]
For example,
\[
\begin{array}{c|l}
    s & \text{basis states in }V_{\mathrm{spin}}\\ \hline
    0 & |0\rangle\\[3pt]
    \frac{1}{2} & |\frac{1}{2}\rangle,\ |{-\frac{1}{2}}\rangle
        \quad\text{also written }|\uparrow\rangle,\ |\downarrow\rangle
        \quad\text{(up, down)}\\[3pt]
    1 & |1\rangle,\ |0\rangle,\ |{-1}\rangle
\end{array}
\]

The existence of spin states is revealed, for example, by the
\underline{Stern--Gerlach experiment}. A magnet produces an inhomogeneous
magnetic field, which exerts different forces on different spin states.
For $s=\frac{1}{2}$, the beam splits into two:
\begin{center}
\setlength{\unitlength}{1mm}
\begin{picture}(135,49)
    \put(5,25){\vector(1,0){26}}
    \put(1,17){\makebox(28,4){beam of particles}}
    \put(36,34){\framebox(25,8){magnet}}
    \put(36,8){\framebox(25,8){magnet}}
    \put(31,25){\line(1,0){25}}
    \qbezier(56,25)(69,25)(88,40)
    \qbezier(56,25)(69,25)(88,10)
    \put(88,40){\vector(4,3){8}}
    \put(88,10){\vector(4,-3){8}}
    \put(100,40){\makebox(28,4)[l]{$|\uparrow\rangle$}}
    \put(100,3){\makebox(28,4)[l]{$|\downarrow\rangle$}}
    \put(20,0){\makebox(66,4){inhomogeneous magnetic field}}
\end{picture}
\end{center}

A general state $|\psi\rangle\in V$ has a $(2s+1)$-component
wavefunction
\[
    \psi(\mathbf{x},r)=\langle\mathbf{x},r|\psi\rangle.
\]
Here $|\psi(\mathbf{x},r)|^2$ is the probability density for the
particle's position with spin state $r$. A mixed or hybrid notation
is also common:
\[
    \sum_r\psi(\mathbf{x},r)|r\rangle,
\]
which is a spin-space-valued function of position.

\subsection*{3.3 Multiparticle: bosons \& fermions}

Consider particles labelled by $a=1,2,\ldots,N$. Let $V_a$ be the
state space for each particle, with basis
$\{|\mathbf{x}_a,r_a\rangle\}$. General multiparticle states belong to
\[
    V_1\otimes V_2\otimes\cdots\otimes V_N,
\]
with basis states
\begin{align*}
    &|\mathbf{x}_1,r_1,\mathbf{x}_2,r_2,\ldots,\mathbf{x}_N,r_N\rangle\\
    &\qquad=|\mathbf{x}_1,r_1\rangle\otimes|\mathbf{x}_2,r_2\rangle
        \otimes\cdots\otimes|\mathbf{x}_N,r_N\rangle.
\end{align*}

\medskip
\noindent
\underline{(a) Identical particles: Bose--Einstein \& Fermi--Dirac statistics}

If the particles are identical, then $V_a\cong V$ for each $a$.
Consider first $N=2$. Define an exchange operator $U$ by its action
on basis states:
\[
    U|\mathbf{x}_1,r_1,\mathbf{x}_2,r_2\rangle
      =|\mathbf{x}_2,r_2,\mathbf{x}_1,r_1\rangle.
\]
For a physical two-particle state,
\[
    U|\Psi\rangle=\eta|\Psi\rangle,
\]
because $|\Psi\rangle$ and $U|\Psi\rangle$ are physically equivalent
states. But $U^2=1$, as follows from its action on basis states, so
\[
    \eta^2=1,\qquad\eta=\pm1.
\]
Thus physical states lie in one of the two subspaces of $V\otimes V$:
\begin{align*}
    (V\otimes V)_{\mathrm{S}}
        &=\operatorname{span}\bigl\{
          |\psi\rangle\otimes|\phi\rangle
          +|\phi\rangle\otimes|\psi\rangle\bigr\},
        &&\eta=+1,\\
    (V\otimes V)_{\mathrm{A}}
        &=\operatorname{span}\bigl\{
          |\psi\rangle\otimes|\phi\rangle
          -|\phi\rangle\otimes|\psi\rangle\bigr\},
        &&\eta=-1,
\end{align*}
where $|\psi\rangle,|\phi\rangle$ range over $V$. These are the
\underline{symmetric} and \underline{antisymmetric} subspaces,
respectively; the displayed spanning states need not be normalised.

Similarly, for $N\geq2$, define $U_{(a,b)}$ to exchange the labels
$(\mathbf{x}_a,r_a)$ and $(\mathbf{x}_b,r_b)$ in a basis state.
For a physical $N$-particle state,
\[
    U_{(a,b)}|\Psi\rangle=\eta_{(a,b)}|\Psi\rangle,
    \qquad\eta_{(a,b)}=\pm1,
\]
since $U_{(a,b)}^2=1$.

For any permutation $\pi$ of $\{1,2,\ldots,N\}$, define
\begin{align*}
    &U_\pi|\mathbf{x}_1,r_1,\mathbf{x}_2,r_2,\ldots,\mathbf{x}_N,r_N\rangle\\
    &\qquad=|\mathbf{x}_{\pi(1)},r_{\pi(1)},
        \mathbf{x}_{\pi(2)},r_{\pi(2)},\ldots,
        \mathbf{x}_{\pi(N)},r_{\pi(N)}\rangle.
\end{align*}
On a physical state,
\[
    U_\pi|\Psi\rangle=\eta_\pi|\Psi\rangle.
\]
The algebra of swaps (transpositions) requires $\eta_{(a,b)}$ to
have the same value for every pair $(a,b)$. Every permutation can
be expressed as a sequence of swaps, so there are two possibilities:
\[
    \eta_\pi=
    \begin{cases}
        1 & \text{for all }\pi\quad\text{(bosons)},\\[3pt]
        \operatorname{sgn}(\pi)=(-1)^{\text{number of swaps in }\pi}
          & \text{for all }\pi\quad\text{(fermions)}.
    \end{cases}
\]
The same alternative must be chosen for all permutations.

\begin{center}
    \textbf{End of lecture 9}
\end{center}
\end{document}
